\chapter{Ajustements de convexité et produits à maturité constante}\label{ch:rc:convexity-adjustments-and-constant-maturity-products}

Une note structurée verse, une fois par an, le taux swap dix ans en euros observé
ce jour-là. Pour le coupon dans dix ans, le taux swap forward dix ans sur les
courbes du chapitre~2 vaut 3.05~\%, et un trader qui coterait le coupon au
forward perdrait de l'argent~: sa juste valeur est de 3.37~\%. Le forward n'a
rien d'erroné. C'est le taux équitable d'un swap qui démarre dans dix ans, une espérance
sous la probabilité dont le numéraire est l'annuité de ce swap~; la note verse le
taux qu'une fois, un an après son observation, et sous la probabilité de cette
date de paiement, l'espérance du taux est plus élevée, d'un montant qui dépend de
la volatilité du taux et de la façon dont l'annuité évolue avec lui. La
différence est un ajustement de convexité. Le livre~2 en a rencontré un sur les
contrats à terme (livre~2, chapitre~8)~; ce chapitre les calcule pour des taux versés
au mauvais moment, dans la mauvaise devise ou avec la mauvaise maturité, les pricer
par réplication à l'aide des swaptions du cube du chapitre~5, et valorise une note sur
l'écart entre deux taux swap.

\section{Verser un taux au mauvais moment~: ajustements de calendrier}

Un taux forward $F_k$ sur $[T_{k-1},T_k]$ est une martingale sous la
probabilité forward $T_k$ (voir la \cref{prop:rc:vanilla-rates-options:measures}), dont
le numéraire correspond à sa date de paiement naturelle. Si on le verse à une
autre date, son espérance change.

\begin{definition}[Ajustement de calendrier]\label{def:rc:convexity-adjustments-and-constant-maturity-products:timing}
L'\emph{ajustement de calendrier}\index{ajustement de calendrier} (timing adjustment) d'un taux versé à une date
$T_p$ autre que sa date de paiement naturelle vaut
$\E^{T_p}[F_k(T_{k-1})]-F_k(0)$~: l'écart entre son espérance sous la
probabilité forward $T_p$ et son forward.
\end{definition}

\begin{proposition}[Un taux versé au moment de sa fixation]\label{prop:rc:convexity-adjustments-and-constant-maturity-products:timing}
Si $F_k$ est versé à sa date de fixation $T_{k-1}$ au lieu de $T_k$,
\[
\E^{T_{k-1}}\bigl[F_k(T_{k-1})\bigr] = F_k(0) +
\frac{\delta_k\,\Var^{T_k}\bigl(F_k(T_{k-1})\bigr)}{1+\delta_kF_k(0)} ,
\]
qui vaut $F_k(0)+\delta_k\sigma_N^2T_{k-1}/(1+\delta_kF_k(0))$ dans le modèle normal,
et $F_k(0)+\delta_kF_k(0)^2\sigma^2T_{k-1}/(1+\delta_kF_k(0))$ au premier ordre dans
le modèle lognormal.
\end{proposition}
\begin{proof}
La densité de $\mathbb Q^{T_{k-1}}$ par rapport à $\mathbb Q^{T_k}$ sur $\mathcal
F_{T_{k-1}}$ est le rapport des numéraires normalisés en zéro,
$\frac{P(T_{k-1},T_{k-1})/P(T_{k-1},T_k)}{P(0,T_{k-1})/P(0,T_k)} =
\frac{1+\delta_kF_k(T_{k-1})}{1+\delta_kF_k(0)}$ (livre~4, chapitre~5).
Donc $\E^{T_{k-1}}[F_k] = \E^{T_k}[F_k(1+\delta_kF_k)]/(1+\delta_kF_k(0))$, et
$\E^{T_k}[F_k^2] = F_k(0)^2+\Var^{T_k}(F_k)$.
\end{proof}

\begin{omfigure}
\begin{tikzpicture}[x=1.6cm,font=\small]
\draw[->,thick] (0,0) -- (6.4,0) node[right] {temps};
\foreach \x/\l in {0/0,3/$T_{k-1}$,5/$T_k$} {\draw (\x,-0.1) -- (\x,0.1); \node[below] at (\x,-0.12) {\l};}
\draw[omDef,thick,fill=omDef!12] (3,0.45) rectangle (5,0.85); \node[font=\footnotesize] at (4,0.65) {accumulation de $F_k$};
\draw[-{Stealth[length=2mm]},omDef,thick] (5,1.6) -- (5,0.95);
\node[text=omDef,align=center,font=\footnotesize] at (5,2.0) {paiement naturel~:\\ $\E^{T_k}[F_k]=F_k(0)$};
\draw[-{Stealth[length=2mm]},omThm,thick] (3,1.6) -- (3,0.95);
\node[text=omThm,align=center,font=\footnotesize] at (3,2.0) {versé à la fixation~:\\ un ajustement de calendrier};
\draw[-{Stealth[length=2mm]},omProp,thick] (5.9,1.6) -- (5.9,0.1);
\node[text=omProp,align=left,font=\footnotesize,anchor=west] at (6.0,1.3) {un taux swap versé une fois,\\ plus tard~: un ajustement CMS};
\end{tikzpicture}

\omcaption{Le moment du versement détermine la probabilité de l'espérance du taux. Versé
à la fin de sa période d'accumulation, un taux forward ne demande aucun ajustement~; versé
à sa fixation, il demande un \omterm{def:rc:convexity-adjustments-and-constant-maturity-products:timing}{ajustement de calendrier}~; un taux swap versé une seule fois, à
une date unique, demande un ajustement de maturité constante.}\label{fig:rc:convexity-adjustments-and-constant-maturity-products:timing}
\end{omfigure}

\begin{example}[Un taux à six mois versé par anticipation]\label{ex:rc:convexity-adjustments-and-constant-maturity-products:timing}
Un taux à six mois dont le forward vaut 3.~\% se fixe dans cinq ans, avec une
volatilité normale de 75 points de base (une \omterm{def:rc:vanilla-rates-options:lognormal}{volatilité lognormale} de 25~\%), versé à
sa fixation~: $0.5\times0.0075^2\times5/(1+0.5\times0.03) = 1.39$ points de base. Petit
pour une période, il est versé à chaque période d'un swap à terme échu et croît
avec la date de fixation.
\end{example}

\section{Swaps à maturité constante et leur réplication par swaptions}

\begin{definition}[Swap à maturité constante]\label{def:rc:convexity-adjustments-and-constant-maturity-products:cms}
Un \emph{swap à maturité constante}\index{swap à maturité constante} (constant-maturity swap, CMS) échange
une jambe variable dont le taux à chaque remise à jour est le taux de swap au pair
d'une maturité fixe (le «~taux CMS~», par exemple le taux swap dix ans observé à la remise à jour)
contre une jambe fixe ou une autre jambe variable. Un caplet CMS verse
$\max(S(T)-K,0)$ sur le taux CMS observé en $T$, une fois, à une date $T_p$.
\end{definition}

\begin{definition}[Ajustement de convexité CMS]\label{def:rc:convexity-adjustments-and-constant-maturity-products:cmsadj}
L'\emph{ajustement de convexité CMS}\index{ajustement de convexité CMS} est
$\E^{T_p}[S_{a,b}(T)]-S_{a,b}(0)$~: l'espérance du taux swap sous la
probabilité forward de sa date de paiement, moins son forward, l'espérance étant
prise sous la probabilité d'annuité.
\end{definition}

Changeons de probabilité, de la probabilité d'annuité vers la date de paiement~:
\[
\E^{T_p}[S(T)] = \frac{A(0)}{P(0,T_p)}\,\E^{a,b}\Bigl[S(T)\,\frac{P(T,T_p)}{A(T)}\Bigr].
\]
Le rapport $P(T,T_p)/A(T)$ est une fonction de toute la courbe en $T$~; pour pricer, il
est ramené à une fonction du seul taux swap.

\begin{definition}[Modèle linéaire du taux swap terminal]\label{def:rc:convexity-adjustments-and-constant-maturity-products:tsr}
Le \emph{modèle linéaire du taux swap terminal}\index{modèle linéaire du taux swap terminal}
pose $P(T,T_p)/A(T) = a_0 + a_1S(T)$, avec $a_0 = 1/\sum_i\delta_i$ (la valeur lorsque
tous les taux sont nuls) et $a_1$ choisi de sorte que la relation soit vraie aujourd'hui~:
$a_0+a_1S(0) = P(0,T_p)/A(0)$.
\end{definition}

\begin{proposition}[Le taux CMS par réplication]\label{prop:rc:convexity-adjustments-and-constant-maturity-products:cms}
Sous le \omterm{def:rc:convexity-adjustments-and-constant-maturity-products:tsr}{modèle linéaire du taux swap terminal},
\[
\E^{T_p}[S(T)] = S(0) + \frac{a_1\,\Var^{a,b}(S(T))}{a_0+a_1S(0)},\qquad
\Var^{a,b}(S(T)) = 2\int_{-\infty}^{\infty} O(K)\,dK,
\]
où $O(K)$ est le prix de la swaption hors de la monnaie par unité d'annuité au strike
$K$ (une receveuse sous le forward, une payeuse au-dessus). Avec un smile normal plat,
$\Var^{a,b}(S(T)) = \sigma_N^2T$ (Hagan, 2003).
\end{proposition}
\begin{proof}
$\E^{a,b}[S(a_0+a_1S)] = a_0S(0) + a_1\bigl(S(0)^2+\Var^{a,b}(S)\bigr)$ puisque $S$ est
une martingale sous $\mathbb Q^{a,b}$~; on divise par $a_0+a_1S(0) = P(0,T_p)/A(0)$. Pour
la variance, $S^2 = S(0)^2 + 2S(0)(S-S(0)) + 2\int_{-\infty}^{S(0)}(K-S)^+dK +
2\int_{S(0)}^{\infty}(S-K)^+dK$ (la réplication statique du livre~5,
chapitre~14), puis on prend les espérances.
\end{proof}

\omcode{code/firm/cms/firm_cms.py}{63}{79}{Le taux CMS par réplication~: swaptions hors de la monnaie intégrées sur le smile.}\label{lst:rc:convexity-adjustments-and-constant-maturity-products:cms}

L'ajustement CMS est donc un prix~: un panier de swaptions hors de la monnaie,
pondéré par $2a_1$. Il est payé par le smile. Avec le smile du cube, les ailes
sont plus épaisses qu'avec une \omterm{def:rc:vanilla-rates-options:flat}{volatilité plate}, la variance est plus élevée, et il en va de même de l'ajustement
(la \cref{fig:rc:convexity-adjustments-and-constant-maturity-products:adj}).

\begin{example}[Swap CMS à dix ans sur le cube euro]\label{ex:rc:convexity-adjustments-and-constant-maturity-products:cmsadj}
Pour le taux swap à dix ans observé dans un, deux, trois, cinq, sept et dix
ans et payé un an plus tard, les ajustements sur les courbes du chapitre~2 valent
1.39, 3.49, 5.87, 11.73, 17.40 et 26.58 points de base avec la volatilité à la
monnaie de chaque section maintenue plate, et 1.62, 4.46, 7.92, 15.51, 22.91 et 32.14
avec le smile du chapitre~5. Le taux CMS à dix ans observé dans dix ans a un
forward de 3.051~\% et une espérance de 3.372~\%~: c'est le coupon du hook.
\end{example}

\begin{omfigure}
\begin{tikzpicture}
\begin{axis}[width=11.5cm,height=5.4cm,
  xlabel={date d'observation du taux swap à dix ans (années)},ylabel={ajustement CMS (pb)},
  tick label style={font=\small},label style={font=\small},
  xmin=0,xmax=10.5,ymin=0,ymax=35,grid=major,grid style={gray!25},
  legend columns=2,legend style={font=\footnotesize,at={(0.5,-0.25)},anchor=north,draw=none,fill=none,/tikz/every even column/.append style={column sep=8pt}},
  table/col sep=comma]
\addplot[omDef,thick,mark=*] table[x=e,y=flat]{figdata/rates-credit-risk/06-convexity-adjustments-and-constant-maturity-products/cmsadj.csv};
\addplot[omThm,very thick,mark=square*] table[x=e,y=smile]{figdata/rates-credit-risk/06-convexity-adjustments-and-constant-maturity-products/cmsadj.csv};
\legend{volatilité normale plate à la monnaie,le smile SABR du cube}
\end{axis}
\end{tikzpicture}

\omcaption{Convexité de l'ajustement du taux CMS euro à dix ans payé un an
après l'observation, par réplication avec une \omterm{def:rc:vanilla-rates-options:flat}{volatilité plate} et avec le smile du
cube du chapitre~5. Le smile ajoute de 16 à 35~\% à l'ajustement. Données~:
courbes du chapitre~2 et cube synthétique du chapitre~5~; le tutoriel du
chapitre.}\label{fig:rc:convexity-adjustments-and-constant-maturity-products:adj}
\end{omfigure}

\begin{omfigure}
\begin{tikzpicture}
\begin{axis}[width=11.5cm,height=5.2cm,
  xlabel={strike de la swaption (\%)},ylabel={pb d'ajustement par 1~\% de strike},
  tick label style={font=\small},label style={font=\small},
  xmin=-3,xmax=7,ymin=0,ymax=9,grid=major,grid style={gray!25},
  table/col sep=comma]
\addplot[omDef,very thick] table[x=k,y=c]{figdata/rates-credit-risk/06-convexity-adjustments-and-constant-maturity-products/contrib.csv};
\draw[gray,dashed] (axis cs:3.051,0) -- (axis cs:3.051,9) node[pos=0.9,right,font=\footnotesize,text=black] {forward 3.05~\%};
\end{axis}
\end{tikzpicture}

\omcaption{D'où vient l'ajustement CMS à dix ans du taux à dix ans~: la
contribution des swaptions de chaque strike, receveuses sous le forward et
payeuses au-dessus, aux 32.14 points de base. Les swaptions situées à un point de pourcentage
ou plus de la monnaie en portent une part importante, ce qui explique l'importance des ailes du
smile. Données~: le tutoriel du chapitre.}\label{fig:rc:convexity-adjustments-and-constant-maturity-products:contrib}
\end{omfigure}

Un caplet CMS se réplique de la même manière. En écrivant $(S-K)^+(a_0+a_1S) =
(a_0+a_1K)(S-K)^+ + a_1((S-K)^+)^2$ et $((S-K)^+)^2 = 2\int_K^\infty(S-L)^+dL$,
sa valeur vaut $A(0)\bigl[(a_0+a_1K)\,\mathrm{Pay}(K) + 2a_1\int_K^\infty
\mathrm{Pay}(L)\,dL\bigr]$, où $\mathrm{Pay}$ est le prix de la payeuse par unité d'annuité.

\omcode{code/firm/cms/firm_cms.py}{82}{90}{Un caplet CMS comme panier de swaptions payeuses à son strike et au-dessus.}\label{lst:rc:convexity-adjustments-and-constant-maturity-products:caplet}

\begin{example}[Un caplet CMS]\label{ex:rc:convexity-adjustments-and-constant-maturity-products:caplet}
Un caplet à 3.50~\% sur le taux à dix ans observé dans cinq ans (forward
3.098~\%, volatilité de smile de 75.9 points de base au strike), payé un an
plus tard, vaut 48.5 points de base de notionnel par réplication~; actualiser un
call de Bachelier sur le forward donnerait 43.0, soit 13~\% de trop peu.
\end{example}

\section{Le modèle linéaire du taux swap terminal en pratique}

Le modèle n'a besoin que de la courbe du jour~: $a_0$ provient des accroissements, $a_1$ de
$P(0,T_p)/A(0)$. Il est exact lorsque la courbe en $T$ est plate et se déplace en
parallèle, et constitue une bonne approximation sinon~; ses alternatives appliquent
$P(T,T_p)/A(T)$ par une formule d'annuité à courbe plate (le modèle standard de Hagan)
ou par les obligations d'un modèle de taux court (chapitre~7). Ce qu'il néglige est la
dépendance de l'application à la forme de la courbe, qui compte pour les longs délais de
paiement et les courbes pentues~; ce qui compte davantage en pratique est le smile loin de
la monnaie, que la réplication permet d'atteindre. Un desk qui valorise des produits CMS
extrapole donc avec soin les ailes de son cube (densités du chapitre~5) et calibre souvent
les paramètres d'ailes sur les spreads CMS négociés plutôt que sur les swaptions.

\section{Ajustements quanto en taux}

Un taux swap euro payé en dollars est un taux observé dans une devise et payé dans
une autre. Sous la probabilité de paiement en dollars, sa dérive change de la covariance
avec le taux de change, l'ajustement quanto du livre~5, chapitre~17.
Pour un taux normal de volatilité $\sigma_N$, une volatilité EURUSD $\sigma_X$
(dollars par euro) et une corrélation $\rho$ entre le taux et l'EURUSD, l'espérance
varie de $-\rho\,\sigma_N\sigma_X\,T$ au premier ordre.

\begin{example}[Un taux euro payé en dollars]\label{ex:rc:convexity-adjustments-and-constant-maturity-products:quanto}
Avec $\sigma_N = 75$ points de base, $\sigma_X = 8~\%$, $\rho = 0.3$ et cinq ans,
l'ajustement vaut $-0.3\times0.0075\times0.08\times5 = -9$ points de base~: quand les
taux euro montent, l'euro tend lui aussi à monter, de sorte que chaque point de base est
payé en euros plus valables que ce que le paiement en dollars restitue.
\end{example}

\section{Options sur spread CMS}

\begin{definition}[Option sur spread CMS]\label{def:rc:convexity-adjustments-and-constant-maturity-products:spread}
Une \emph{option sur spread CMS}\index{option sur spread CMS} verse $\max(S_1(T)-S_2(T)-K,0)$
en $T_p$ sur la différence de deux taux CMS de maturités différentes observés à la
même date, typiquement le dix ans et le deux ans~; un panier de ces options avec
$K=0$ verse un coupon de steepener.
\end{definition}

Évaluée avec chaque taux normal autour de sa propre espérance CMS, le spread est
normal de volatilité $\sigma_s = \sqrt{\sigma_1^2+\sigma_2^2-2\rho\sigma_1\sigma_2}$~:
le prix dépend de la corrélation entre les deux taux, qu'aucune swaption ne cote. Le
desk qui vend une note steepener est vendeur de volatilité du spread et donc acheteur de
corrélation~; sa couverture repose sur des produits sensibles à la corrélation, et sa
valorisation dépend d'une estimation.

\begin{example}[La volatilité du spread]\label{ex:rc:convexity-adjustments-and-constant-maturity-products:spreadvol}
Sur 2016--2026, les variations quotidiennes des taux au pair du Trésor américain à deux et dix ans
ont présenté une corrélation de 0.769 (données du chapitre~3, en guise de proxy des taux swap). Avec les
volatilités à la monnaie du cube pour l'échéance d'un an, le spread des
taux euro à dix et deux ans affiche une volatilité normale de 45.1 points de base, nettement
inférieure à celle de chacun des deux taux.
\end{example}

\begin{omfigure}
\begin{tikzpicture}
\begin{axis}[width=11cm,height=5.2cm,
  xlabel={corrélation des taux à dix et deux ans},ylabel={participation juste},
  tick label style={font=\small},label style={font=\small},
  xmin=0.48,xmax=1.0,ymin=2,ymax=5.4,grid=major,grid style={gray!25},
  table/col sep=comma]
\addplot[omThm,very thick,mark=*,mark size=1.5pt] table[x=rho,y=p]{figdata/rates-credit-risk/06-convexity-adjustments-and-constant-maturity-products/steep.csv};
\draw[gray,dashed] (axis cs:0.769,2) -- (axis cs:0.769,5.4) node[pos=0.9,left,font=\footnotesize,text=black] {Estimation du Trésor 0.769};
\end{axis}
\end{tikzpicture}

\omcaption{Participation juste d'une note steepener à dix ans (le problème du
week-end) en fonction de la corrélation entre les taux swap à dix et deux ans~:
plus la corrélation est élevée, moins le spread est volatil et moins chaque unité de son floor
coûte cher, si bien que la note peut en verser davantage. Données~:
tutoriel du chapitre.}\label{fig:rc:convexity-adjustments-and-constant-maturity-products:steep}
\end{omfigure}

\section{Tutoriel~: réplication d'un taux CMS}\label{tut:rc:convexity-adjustments-and-constant-maturity-products:tutorial}

\begin{tutorial}
\textbf{Objectif.} Calculer les ajustements CMS par réplication sur le cube, pricer un caplet CMS
et une note steepener. \textbf{État final~:}
\cref{fig:rc:convexity-adjustments-and-constant-maturity-products:adj,fig:rc:convexity-adjustments-and-constant-maturity-products:contrib,fig:rc:convexity-adjustments-and-constant-maturity-products:steep}.
\begin{tutsteps}
\item \textbf{Mise en place}~: \texttt{setup(e, tenor)} renvoie le taux swap forward, son
annuité, le facteur d'actualisation du paiement et les accroissements issus des courbes du chapitre~2.
\item \textbf{Ajustements}~: \texttt{cms\_table()} avec un smile plat
(\texttt{cms\_rate\_hagan}) et avec le cube (\texttt{cms\_rate})~; le modèle décalé du cube
n'a aucune masse en dessous de $-3~\%$, si bien que l'intégrale part de ce niveau.
\item \textbf{D'où il provient}~: \texttt{replication\_contributions()}.
\item \textbf{La note}~: \texttt{steepener()}~; \texttt{fig\_rc\_cms.py} écrit les
graphiques.
\end{tutsteps}
\textbf{Pour aller plus loin.} Remplacer le smile du cube par sa \omterm{def:rc:vanilla-rates-options:flat}{volatilité plate} à la monnaie
dans le steepener et comparer~; déplacer la date de paiement à la fin du
swap sous-jacent et observer la croissance de l'ajustement.
\end{tutorial}

\section{Construction~: pricing CMS}\label{bld:rc:convexity-adjustments-and-constant-maturity-products:cms}

\begin{build}
\bfield{Rôle} Pricing de tout produit versant un taux au mauvais moment, dans la
mauvaise devise ou avec une maturité constante~; les jambes CMS des
notes structurées du livre.
\bfield{Interface} \texttt{timing\_adjustment\_black},
\texttt{timing\_adjustment\_normal}~; \texttt{linear\_tsr(annuity0, df\_pay, fwd,
accruals)}~; \texttt{otm\_integral}, \texttt{cms\_rate}, \texttt{cms\_caplet},
\texttt{cms\_rate\_hagan}~; \texttt{quanto\_adjustment\_normal}~;
\texttt{spread\_vol}, \texttt{cms\_spread\_option}.
\bfield{Règles} Le smile est une fonction quelconque qui va du strike à la volatilité normale (le
cube du chapitre~5 du livre)~; l'intégrale est bornée inférieurement par le plancher du modèle de smile~;
numpy seulement.
\bfield{Tests d'acceptation} \texttt{code/firm/cms/tests/}~: avec un smile plat, l'ajustement
répliqué égale la formule fermée~; sans pente dans l'application annuité, il
n'y a aucun ajustement~; un caplet CMS profondément dans la monnaie vaut le
forward CMS actualisé~; les \omterm{def:rc:convexity-adjustments-and-constant-maturity-products:timing}{ajustements de calendrier} concordent entre les formes de Black et normale
à volatilités appariées~; une corrélation parfaite supprime la valeur temps de l'option sur
spread.
\bfield{Pour aller plus loin} Application de Hagan au taux standard (annuité à courbe plate)~; \omterm{def:rc:convexity-adjustments-and-constant-maturity-products:spread}{options sur spread CMS}
avec une copule des deux smiles~; la jambe complète d'un swap CMS avec retards de paiement.
\end{build}

\begin{omsources}
\item P.\ S.\ Hagan, ``Convexity conundrums: pricing CMS swaps, caps, and
floors'', \emph{Wilmott Magazine}, mars 2003.
\item L.\ Andersen et V.\ Piterbarg, \emph{Interest Rate Modeling}, volume 3,
Atlantic Financial Press, 2010, chapitres 16--17.
\item US Department of the Treasury, \emph{Daily Treasury Par Yield Curve Rates}
(l'estimation de la corrélation).
\end{omsources}

\section{Exercices}

\begin{exercise}[$\star$]\label{exo:rc:convexity-adjustments-and-constant-maturity-products:1}
Un taux à trois mois de 2.5~\% fixé dans quatre ans, avec 80 points de base de
volatilité normale, est payé à sa fixation. Calculer son \omterm{def:rc:convexity-adjustments-and-constant-maturity-products:timing}{ajustement de calendrier}.
\end{exercise}

\begin{exercise}[$\star$]\label{exo:rc:convexity-adjustments-and-constant-maturity-products:2}
Pourquoi l'ajustement CMS est-il positif pour un taux payé peu après son observation~?
Quel est le signe de $a_1$ lorsque les taux sont positifs~?
\end{exercise}

\begin{exercise}[$\star$]\label{exo:rc:convexity-adjustments-and-constant-maturity-products:3}
Pour un swap annuel à dix ans avec des accroissements de 1.0139 chacun, calculer $a_0$.
\end{exercise}

\begin{exercise}[$\star\star$]\label{exo:rc:convexity-adjustments-and-constant-maturity-products:4}
Avec un smile plat, l'ajustement à cinq ans de l'\cref{ex:rc:convexity-adjustments-and-constant-maturity-products:cmsadj}
vaut 11.73 points de base. Quelle serait son ampleur à une volatilité normale supérieure de 20~\%~?
\end{exercise}

\begin{exercise}[$\star\star$]\label{exo:rc:convexity-adjustments-and-constant-maturity-products:5}
Lire la \cref{fig:rc:convexity-adjustments-and-constant-maturity-products:contrib}~:
environ quelle part de l'ajustement provient des swaptions situées à plus de deux
points de pourcentage du forward~? Qu'est-ce que cela indique sur la manière dont les
ailes du cube sont valorisées~?
\end{exercise}

\begin{exercise}[$\star\star$]\label{exo:rc:convexity-adjustments-and-constant-maturity-products:6}
Dans l'\cref{ex:rc:convexity-adjustments-and-constant-maturity-products:quanto}, que devient
l'ajustement si la corrélation vaut $-0.3$, et pourquoi~?
\end{exercise}

\begin{exercise}[$\star\star\star$]\label{exo:rc:convexity-adjustments-and-constant-maturity-products:7}
\emph{Programmation.} Calculer l'\omterm{def:rc:convexity-adjustments-and-constant-maturity-products:cmsadj}{ajustement de convexité CMS} à dix ans du taux observé dans cinq
ans, lorsqu'il est payé à la fin du swap sous-jacent (quinze ans) plutôt
qu'un an après l'observation. Expliquer le signe.
\end{exercise}

\begin{exercise}[$\star\star\star$]\label{exo:rc:convexity-adjustments-and-constant-maturity-products:8}
\emph{Trouver la faille.} «~Un swap CMS est linéaire en les taux, donc on valorise ses coupons aux
taux swap forward et on les couvre avec des swaps à départ différé.~» Le corriger.
\end{exercise}

\section{Problème~: La note steepener}

\begin{problem}[{Problème du week-end --- évaluation d'un steepener à dix ans}]\label{pb:rc:convexity-adjustments-and-constant-maturity-products:1}
Une banque émet une note en euros à dix ans qui verse chaque année $k=1,\dots,10$ un coupon de
$p\times\max(S_{10}(k)-S_2(k),0)$, soit le taux swap à dix ans moins le taux swap à deux ans
observé à l'année $k$, payé à $k+1$. L'objectif de financement de la banque est de verser des coupons
d'un montant équivalent à celui d'un coupon fixe de 2~\% aux mêmes dates. Les taux sont ceux du
chapitre~2, les smiles sont ceux du chapitre~5, et la corrélation entre les deux taux
est estimée à 0.769.

\medskip\noindent\textbf{Partie I --- Les composantes.}
\begin{enumerate}
\item Qu'est-ce que chaque coupon, vu comme une option~?
\item Pourquoi chaque taux a-t-il besoin d'un \omterm{def:rc:convexity-adjustments-and-constant-maturity-products:cmsadj}{ajustement de convexité CMS}, et lequel est le plus important~?
\item Donner la valeur de la jambe fixe à 2~\% par unité de notionnel.
\item Donner la valeur de la jambe de coupons avec $p=1$.
\item Donner la participation juste $p$.
\end{enumerate}

\medskip\noindent\textbf{Partie II --- Corrélation.}
\begin{enumerate}[resume]
\item Donner la participation juste pour des corrélations de 0.6, 0.8 et 0.95.
\item Pourquoi augmente-t-elle avec la corrélation~?
\item Dans quel sens la banque émettrice est-elle exposée à la corrélation après avoir vendu la note~?
\item Donner la volatilité normale du spread de la première année.
\item Pourquoi les données Treasury ne sont-elles ici qu'un proxy~?
\end{enumerate}

\medskip\noindent\textbf{Partie III --- Risque.}
\begin{enumerate}[resume]
\item Quelles swaptions couvrent les \omterm{def:rc:convexity-adjustments-and-constant-maturity-products:cmsadj}{ajustements de convexité CMS}~?
\item Que se passe-t-il pour la valeur de la note si les ailes du smile montent~?
\item La valeur de la note est-elle sensible au niveau des taux, à la pente, ou aux deux~?
\item Que se passe-t-il pour l'investisseur si la courbe s'inverse~?
\item Pourquoi les portefeuilles des dealers sur de telles notes deviennent-ils importants et déséquilibrés~?
\end{enumerate}

\medskip\noindent\textbf{Partie IV --- Jugement.}
\begin{enumerate}[resume]
\item Quelle participation coter, et pourquoi pas 2.96~?
\item Comment valoriser la corrélation en l'absence de cotations~?
\item Que répondre à un investisseur qui affirme que la courbe «~se repentifie toujours~» de nouveau~?
\item Énoncer le \emph{résultat nommé}~: la participation juste au niveau de la corrélation estimée,
et son intervalle entre des corrélations de 0.6 et 0.95.
\item En une phrase~: qu'achète un investisseur en steepener~?
\end{enumerate}
\end{problem}

\section{Questions d'entretien}

\begin{interviewq}[$\star$ \iqroles{chercheur, trader}]\label{iq:rc:convexity-adjustments-and-constant-maturity-products:1}
Qu'est-ce qu'un ajustement de convexité, en une phrase, et en donner deux exemples~?
\end{interviewq}

\begin{interviewq}[$\star\star$ \iqroles{chercheur}]\label{iq:rc:convexity-adjustments-and-constant-maturity-products:2}
Dériver l'\omterm{def:rc:convexity-adjustments-and-constant-maturity-products:timing}{ajustement de calendrier} d'un taux payé à sa date de fixing.
\end{interviewq}

\begin{interviewq}[$\star\star$ \iqroles{chercheur, banque}]\label{iq:rc:convexity-adjustments-and-constant-maturity-products:3}
Comment valorisez-vous un caplet CMS à partir de swaptions~?
\end{interviewq}

\begin{interviewq}[$\star\star$ \iqroles{trader}]\label{iq:rc:convexity-adjustments-and-constant-maturity-products:4}
Vous êtes vendeur de coupons CMS. Dans quel sens souhaitez-vous que la volatilité implicite évolue, et
quels strikes comptent le plus~?
\end{interviewq}

\begin{interviewq}[$\star\star\star$ \iqroles{chercheur, risque}]\label{iq:rc:convexity-adjustments-and-constant-maturity-products:5}
Une \omterm{def:rc:convexity-adjustments-and-constant-maturity-products:spread}{option sur spread CMS} dépend d'une corrélation qui n'est pas cotée. Comment la valorisez-vous
et la couvrez-vous~?
\end{interviewq}

\begin{interviewq}[$\star\star\star$ \iqroles{développeur}]\label{iq:rc:convexity-adjustments-and-constant-maturity-products:6}
Votre réplication CMS donne un résultat différent chaque fois que la grille de strikes
change. Que vérifiez-vous~?
\end{interviewq}
